% =============================================================================
% SECTION 4: NUMERICAL DISCRETIZATION
% =============================================================================

\section{Numerical Discretization Approach}

\subsection{Finite Difference Method}

The beam BVP~\eqref{eq:eb-prismatic} is discretized using the \textbf{finite difference method (FDM)}.
The domain $[0, L]$ is partitioned into $N$ equal sub-intervals of width $h = L/N$, yielding $N+1$ grid points:
\[
  x_i = i\,h, \quad i = 0, 1, \ldots, N
\]

The fourth-order derivative is approximated using the standard second-order central difference stencil:

\begin{equation}
  \frac{d^4 w}{dx^4}\bigg|_{x_i}
  \approx
  \frac{w_{i-2} - 4w_{i-1} + 6w_i - 4w_{i+1} + w_{i+2}}{h^4}
  + \mathcal{O}(h^2)
  \label{eq:fd-stencil}
\end{equation}

The discretized governing equation at each \emph{interior} point $i = 2, \ldots, N-2$ is:

\begin{equation}
  \frac{EI}{h^4}\Bigl(w_{i-2} - 4w_{i-1} + 6w_i - 4w_{i+1} + w_{i+2}\Bigr) = q_i
  \label{eq:fd-system}
\end{equation}

where $q_i = q(x_i)$ is the applied load at node $i$.

\subsection{Assembly into a Linear System}

Equation~\eqref{eq:fd-system} applied at all interior nodes yields a banded linear system:

\begin{equation}
  \mathbf{K}\,\mathbf{w} = \mathbf{f}
  \label{eq:linear-system}
\end{equation}

where $\mathbf{K} \in \mathbb{R}^{(N+1)\times(N+1)}$ is the stiffness matrix (pentadiagonal, bandwidth 5), $\mathbf{w}$ is the vector of unknown deflections, and $\mathbf{f}$ contains the load values scaled by $h^4/(EI)$.

\textbf{Simply-supported boundary conditions} are enforced as follows.
The Dirichlet conditions $w_0 = 0$ and $w_N = 0$ are imposed directly on rows 0 and $N$:
\begin{align}
  w_0 = 0 \quad &\Rightarrow \quad \text{row } 0: \; [1,\, 0,\, 0,\, \ldots] \\
  w_N = 0 \quad &\Rightarrow \quad \text{row } N: \; [\ldots,\, 0,\, 0,\, 1]
\end{align}

The zero-moment conditions $w''(0) = 0$ and $w''(L) = 0$ are handled via \textbf{ghost points}.
Applying the centered second-difference at each boundary with $w_0 = w_N = 0$:
\begin{align}
  \frac{w_{-1} - 2w_0 + w_1}{h^2} = 0 \;\wedge\; w_0 = 0
    \quad &\Rightarrow \quad w_{-1} = -w_1 \\
  \frac{w_{N-1} - 2w_N + w_{N+1}}{h^2} = 0 \;\wedge\; w_N = 0
    \quad &\Rightarrow \quad w_{N+1} = -w_{N-1}
\end{align}

These ghost-point values are substituted into the 5-point stencil at the \textbf{near-boundary nodes} $i = 1$ and $i = N-1$, yielding modified equations for rows 1 and $N-1$:
\begin{align}
  i = 1: &\quad \frac{EI}{h^4}\bigl(5w_1 - 4w_2 + w_3\bigr) = q_1
    \label{eq:fd-row1} \\
  i = N-1: &\quad \frac{EI}{h^4}\bigl(w_{N-3} - 4w_{N-2} + 5w_{N-1}\bigr) = q_{N-1}
    \label{eq:fd-rowNm1}
\end{align}
These modified stencils maintain $\mathcal{O}(h^2)$ consistency since the ghost-point approximation introduces only $\mathcal{O}(h^2)$ error in the moment condition.

The complete row structure of $\mathbf{K}$ is therefore:
\begin{itemize}
  \item Row 0: enforces $w_0 = 0$ (Dirichlet)
  \item Row 1: modified 3-term stencil from ghost-point substitution, Eq.~\eqref{eq:fd-row1}
  \item Rows 2 to $N-2$: standard 5-point stencil, Eq.~\eqref{eq:fd-system}
  \item Row $N-1$: modified 3-term stencil from ghost-point substitution, Eq.~\eqref{eq:fd-rowNm1}
  \item Row $N$: enforces $w_N = 0$ (Dirichlet)
\end{itemize}

The system~\eqref{eq:linear-system} is solved using a direct banded solver (LU factorization), exploiting the half-bandwidth-2 (5-diagonal) structure for $\mathcal{O}(N)$ cost.

\subsection{Grid Convergence}

The formal truncation error of stencil~\eqref{eq:fd-stencil} is second-order: $\tau = \mathcal{O}(h^2)$.
Therefore, as the mesh is refined ($N \to \infty$, $h \to 0$), the numerical solution converges to the exact solution at rate $p = 2$ in the $L^\infty$ norm:
\begin{equation}
  \|w_h - w_{\text{exact}}\|_\infty \propto h^2
  \label{eq:convergence}
\end{equation}

A grid convergence study will be performed in the code verification homework by computing the observed order of accuracy $\hat{p}$ from solutions on systematically refined grids ($h$, $h/2$, $h/4$) and comparing against the theoretical value $p = 2$.

\subsection{Uncertainty Quantification}

For uncertainty propagation, the linear solve is repeated $N_{\text{MC}}$ times with independently sampled load and material parameters:
\begin{align}
  T_w &\sim \text{Uniform}(T_{w,\min},\, T_{w,\max}) \quad &\text{(geometric ambiguity)} \\
  q_{\text{roof}} &\sim \mathcal{N}(\mu_{q,\text{roof}},\, \sigma_{q,\text{roof}}^2) \quad &\text{(code + variability)} \\
  q_{\text{floor}} &\sim \mathcal{N}(\mu_{q,\text{floor}},\, \sigma_{q,\text{floor}}^2) \\
  E &\sim \mathcal{N}(\mu_E,\, \sigma_E^2) \quad &\text{(wood grade variability)}
\end{align}

For each sample $j = 1, \ldots, N_{\text{MC}}$, the beam is solved and the performance criteria are evaluated.
The probability of failure is estimated as:
\begin{equation}
  \hat{P}_{\text{fail}} = \frac{1}{N_{\text{MC}}} \sum_{j=1}^{N_{\text{MC}}}
  \mathbf{1}\!\left(\sigma_{\max}^{(j)} > F_b \;\vee\; \tau_{\max}^{(j)} > F_v \;\vee\; w_{\max}^{(j)} > \frac{L}{240}\right)
  \label{eq:pfail}
\end{equation}

Since each FD solve completes in $\mathcal{O}(N)$ time and typical beam discretizations use $N = 50$--$500$ nodes, running $N_{\text{MC}} = 10{,}000$ Monte Carlo samples takes well under one second on modern hardware.
Also note that $L/240$ in the indicator function must be evaluated in the same units as $w_{\max}$ (inches).

\subsection{Summary}

Table~\ref{tab:numerical-summary} summarizes the numerical approach.

\begin{table}[ht]
\centering
\caption{Summary of numerical discretization choices.}
\label{tab:numerical-summary}
\begin{tabular}{ll}
\toprule
\textbf{Attribute} & \textbf{Choice} \\
\midrule
Governing equation & Euler--Bernoulli beam ODE (4th order, 1D) \\
Discretization method & Finite differences (FDM) \\
Spatial stencil & Central differences, $\mathcal{O}(h^2)$ \\
Linear solver & Banded LU factorization \\
Boundary conditions & Simply-supported (primary); fixed-fixed (secondary) \\
Typical grid size & $N = 100$--$500$ nodes \\
UQ method & Monte Carlo sampling ($N_{\text{MC}} = 10{,}000$) \\
Verification benchmark & Analytical solution, Eq.~\eqref{eq:exact-ss} \\
\bottomrule
\end{tabular}
\end{table}
